\section{Statement, conventions, and conditional inputs}
\label{in:statement}

Campana's theory of special varieties and geometric orbifolds
\cite{Campana2004,CampanaOrbifolds} leads to the orbifold Abelianity
Conjecture \cite[Conjecture~12.10(1)]{CampanaOrbifolds}. This note
gives a conditional proof for a special smooth projective fourfold pair
with one smooth divisor of multiplicity two. It accompanies \emph{A
conditional abelianity theorem for special fourfold pairs with a
half-weight divisor}. The main article uses Theorem~1.1 of the companion
paper \cite{Supplied} as its sole result from that paper. The present proof uses its
Theorems~1.1, 2.1, 2.3, 2.5 and 4.1, stated below, and develops the
descent through inertia and quotient minimality in detail.

All varieties in the proof are over $\mathbb C$.  A fibration is a
dominant map with connected general fibre.  For a rational or meromorphic
fibration, this convention is applied on smooth models after resolution.
All fundamental groups are entire discrete topological fundamental
groups.  Base points and transporting paths are chosen compatibly; a
change of these choices conjugates the subgroup images under discussion.

Let $X$ be smooth and let $D$ be a smooth reduced divisor.  The
\emph{order-two root orbifold} $\mathcal X$ of $D$ agrees with $X$ off
$D$, and along $D=\{z_1=0\}$ it has the chart
\[
 z_1=w_1^2,\qquad z_j=w_j\quad (j>1),
\]
with the group $\mu_2$ acting by sign on $w_1$.  Orbifold line bundles,
forms, and their maps are equivariant holomorphic objects on these
charts, with compatible gluing.  We also use finite integral root
orbifolds on smooth spaces with simple normal crossing divisors.  An
order $m_i$ along $D_i$ corresponds to the coefficient $1-1/m_i$;
order one contributes no boundary.

For an orbifold line bundle $\mathcal L$, its Iitaka dimension is
$-\infty$ if every positive tensor power has zero global section space,
and otherwise is the largest dimension of the images of its complete
plurisection maps.  A root orbifold $\mathcal X$ of dimension $n$ is
\emph{special} if there is no nonzero sheaf map
\[
 \mathcal L\longrightarrow\Omega^p_{\mathcal X},
 \qquad 1\le p\le n,
 \qquad \kappa(\mathcal X,\mathcal L)=p.
\]
For a manifold the same definition uses ordinary line bundles and
ordinary differential forms.  A point is special.

\begin{theorem}\label{in:main}
Assume the five results in
Assumptions~\ref{in:input-manifold}--\ref{in:input-positive}.
Let $X$ be a connected
smooth projective variety of dimension four, and let $D$ be a nonempty
smooth connected reduced effective divisor.  If the order-two root
orbifold $\mathcal X$ of $D$ is special, then
\[
 G=\pi_1(X\setminus D,x)\big/\langle\!\langle\gamma^2\rangle\!\rangle
\]
is virtually abelian.  Here $\gamma$ is a positively oriented meridian
of $D$, transported to $x\in X\setminus D$, and virtual abelianity
means the existence of an abelian subgroup of finite index.
\end{theorem}

After the cases of a trivial meridian or one detected by a finite
quotient, the remaining obstruction is the kernel $K$ of
$G\to\pi_1(X)$. We identify it with the image of
the group of an open Albanese fibre. A first minimum, denoted by
$\ell$, ranges over rational fibrations of suitable compact Albanese
fibres whose open fibre images leave an infinite quotient of $K$.
It provides the dimension bound required for every infinite quotient
on one fixed auxiliary manifold in Theorem~4.1 of the companion paper.

This gives bigness on the general fibres of an intermediate fibration
over the Albanese torus. A second minimum, denoted by $b$, ranges over
the Iitaka dimensions of the adjoints defined by actual meridian
orders, on all prepared models and after all torus isogenies.
It provides a lower bound throughout the subsequent changes of model.
Descent through the Iitaka fibration then gives a big adjoint on a
base of dimension $b$. Its plurisections pull back to a differential
line of Iitaka dimension $b$ on a special root orbifold, a contradiction.
The auxiliary manifolds supply the ordinary compact groups required
by the reduction and positivity inputs. Specialness and the final
contradiction remain on the original root orbifold and its finite
coarse \'etale covers.

We use the following five results from the companion paper \cite{Supplied}.

\begin{assumption}[Companion Theorem 1.1]\label{in:input-manifold}
Every connected smooth compact K\"ahler manifold which is special in
the differential-line sense above has virtually abelian fundamental
group, in every complex dimension.
\end{assumption}

\begin{assumption}[Companion Theorem 2.1]\label{in:input-special}
For smooth compact K\"ahler manifolds, specialness is bimeromorphically
invariant and is preserved by connected finite \'etale covers.  The
Albanese morphism of a special manifold is surjective with connected
fibres.  If a holomorphic line bundle $L$ on a smooth compact K\"ahler
manifold $M$ admits a nonzero map $L\to\Omega_M^p$, then
$\kappa(M,L)\le p$.
\end{assumption}

\begin{assumption}[Companion Theorem 2.3]\label{in:input-reduction}
For a smooth compact K\"ahler manifold $M$ and a homomorphism
$\rho:\pi_1(M)\to P$, there is a meromorphic fibration
$\gamma_\rho:M\dashrightarrow B_\rho$ whose general fibre group has
finite image under $\rho$, and which contracts every compact
irreducible subvariety through a very general point whose fundamental
group has finite image under $\rho$.  These subvariety groups may be
computed on smooth resolutions.  The base has a compact K\"ahler
model.  Write
\[
 \operatorname{gdim}(M,\rho)=\dim B_\rho.
\]
\end{assumption}

Here and below, ``very general'' permits the exclusion of a countable
union of proper analytic subvarieties.  No linearity hypothesis is
imposed on $P$.

\begin{assumption}[Companion Theorem 2.5]\label{in:input-tori}
Let $Z$ be a smooth compact K\"ahler manifold, let $\Theta$ be an
effective rational divisor for which $(Z,\Theta)$ is Kawamata log
terminal (klt), and let
$g:Z\to A$ be a fibration onto a complex torus.  For a general smooth
fibre $F$,
\[
 \kappa(Z,K_Z+\Theta)\ge
 \kappa(F,K_F+\Theta|_F).
\]
If $\kappa(Z,K_Z+\Theta)=0$, the Albanese morphism of $Z$ is
surjective with connected fibres.
\end{assumption}

The torus in this statement may be a point.  In all our uses the spaces
are projective, and the boundaries have simple normal crossing support
and coefficients strictly less than one, which verifies the klt
hypothesis.  We use this formulation from the companion paper directly; see also the
results over tori in \cite{Wang}.

\begin{assumption}[Companion Theorem 4.1]\label{in:input-positive}
Let $Y$ be a connected smooth compact K\"ahler manifold of dimension
$d>0$, let
\[
 \Delta=\sum_i(1-m_i^{-1})D_i,
 \qquad m_i\in\mathbb Z,\quad m_i\ge2,
\]
have simple normal crossing support, and let $\mathcal Y$ be its root
orbifold.  Suppose that
\[
 q:\pi_1^{\mathrm{orb}}(\mathcal Y)\twoheadrightarrow Q
\]
is a surjection, that $Q$ is infinite, and that every finite-dimensional
complex linear representation of $Q$ has finite image.

Suppose a connected smooth compact K\"ahler manifold $S$ admits a
dominant meromorphic map $a:S\dashrightarrow Y$ and a homomorphism
$\nu:\pi_1(S)\to Q$ with finite-index image.  Require compatibility
with $q$ on a dense analytic Zariski-open subset $S^\circ\subset S$
where $a$ is holomorphic with values in $Y\setminus\operatorname{Supp}
\Delta$: the restriction of $\nu$ is the composition of $a_*$ with
the natural map to $\pi_1^{\mathrm{orb}}(\mathcal Y)$ and then $q$.

Finally, for every finite-index subgroup $Q_0\le Q$, suppose there is a
connected finite \'etale cover $S'\to S$ such that, writing $\nu'$ for
the induced homomorphism and $Q'=\operatorname{im}\nu'$, one has
$Q'\subseteq Q_0$ and
\[
 \operatorname{gdim}(S',\psi\circ\nu')\ge d
 \quad\text{for every surjection }\psi:Q'\twoheadrightarrow H
 \text{ with }H\text{ infinite}.
\]
Then $K_Y+\Delta$ is big.
\end{assumption}

The descent below follows the strategy of Theorem~5.8 of the companion paper,
and the first dimension minimum follows the strategy of its Section~6.
Sections~\ref{in:minimality}--\ref{in:descent} establish the required
steps for the present orbifold source, including the compatible
dominating manifold and the full quantified quotient condition of
Assumption~\ref{in:input-positive}. Actual meridian orders determine
$\Delta$ on every model. The established orbifold specialness results
used separately are stated and checked in the next section.
